% =============================================================================
%  Resumen Capítulo 1 — INTRODUCCIÓN
%  Libro: Evsukoff, A. (2020). Inteligência Computacional: Fundamentos e
%         Aplicações. Río de Janeiro: E-papers.
% =============================================================================
% =============================================================================
%  preambulo.tex — Preámbulo de los resúmenes en español
%  Evsukoff, A. (2020). Inteligência Computacional. Rio de Janeiro: E-papers.
% =============================================================================
\documentclass[11pt,a4paper]{article}

\usepackage[T1]{fontenc}
\usepackage[utf8]{inputenc}
\usepackage[spanish,es-tabla]{babel}
\usepackage{lmodern}
\usepackage[a4paper,margin=2.4cm]{geometry}
\usepackage{amsmath,amssymb,amsthm,mathtools}
\usepackage{bm}
\usepackage{enumitem}
\usepackage{xcolor}
\usepackage{tcolorbox}
\tcbuselibrary{skins,breakable}
\usepackage{hyperref}
\usepackage{microtype}
\usepackage{titlesec}
\usepackage{fancyhdr}
\usepackage{booktabs}
\usepackage{array}

\hypersetup{
  colorlinks=true,
  linkcolor=black!70,
  urlcolor=blue!60!black,
  citecolor=black!70,
  pdfborder={0 0 0}
}

\titleformat{\section}{\Large\bfseries\sffamily\color{black!85}}{\thesection}{0.7em}{}
\titleformat{\subsection}{\large\bfseries\sffamily\color{black!75}}{\thesubsection}{0.6em}{}

\pagestyle{fancy}
\fancyhf{}
\fancyhead[L]{\small\itshape Resumen — Inteligencia Computacional (Evsukoff, 2020)}
\fancyhead[R]{\small\thepage}
\renewcommand{\headrulewidth}{0.3pt}

\newtcolorbox{conceito}[1][]{
  colback=blue!4, colframe=blue!55!black, boxrule=0.6pt,
  arc=2mm, left=6pt, right=6pt, top=4pt, bottom=4pt,
  title={\bfseries Concepto clave}, fonttitle=\sffamily\bfseries,
  breakable, #1
}
\newtcolorbox{formula}[1][]{
  colback=gray!6, colframe=gray!50, boxrule=0.4pt,
  arc=1.5mm, left=6pt, right=6pt, top=3pt, bottom=3pt,
  breakable, #1
}
\newtcolorbox{aplicacao}[1][]{
  colback=green!4, colframe=green!45!black, boxrule=0.5pt,
  arc=2mm, left=6pt, right=6pt, top=3pt, bottom=3pt,
  title={\bfseries Aplicación}, fonttitle=\sffamily\bfseries,
  breakable, #1
}

\newcommand{\R}{\mathbb{R}}
\newcommand{\E}{\mathbb{E}}
\newcommand{\Var}{\operatorname{Var}}
\newcommand{\Cov}{\operatorname{Cov}}
\newcommand{\diag}{\operatorname{diag}}
\newcommand{\argmin}{\operatorname*{arg\,min}}
\newcommand{\argmax}{\operatorname*{arg\,max}}
\newcommand{\T}{^{\!\top}}
\DeclareMathOperator{\sign}{sign}
\DeclareMathOperator{\tr}{tr}


\title{\sffamily\bfseries Capítulo 1 — Introducción\\
       \large Inteligencia Computacional para Ciencia de Datos}
\author{Resumen de estudio}
\date{}

\begin{document}
\maketitle
\thispagestyle{fancy}

\begin{conceito}
La \textbf{inteligencia computacional} es un conjunto de paradigmas
computacionales para el desarrollo y aplicación de sistemas inteligentes
en problemas complejos del mundo real. Incluye métodos estadísticos,
aprendizaje de máquina, sistemas \emph{fuzzy}, redes neuronales, máquinas
de vectores de soporte y algoritmos evolutivos, frecuentemente integrados
en abordajes híbridos. La \textbf{ciencia de datos} es la disciplina que
reúne matemáticas, estadística, computación e ingeniería con el
conocimiento del dominio de cada aplicación para extraer modelos y
conocimiento a partir de datos.
\end{conceito}

\section{De la inteligencia artificial a la inteligencia computacional}

La pregunta filosófica ``¿puede pensar una máquina?'', formulada por
Wittgenstein en los años 1920 y operacionalizada por Alan Turing en 1950
con el célebre \textbf{test de Turing}, marca el inicio de la reflexión
sobre el comportamiento inteligente en máquinas. La inteligencia
artificial (IA) nace a fines de los 1950 de la integración entre lógica,
ciencias de la computación y neurofisiología. Tempranamente se distinguen
dos corrientes:

\begin{itemize}[leftmargin=1.2cm]
  \item \textbf{Simbolista}: lógicos matemáticos que construyen
        demostradores de teoremas y, más tarde, \emph{sistemas expertos}
        basados en reglas y hechos.
  \item \textbf{Conexionista}: investigadores que simulan el procesamiento
        cerebral mediante neuronas artificiales. El \emph{Perceptrón} de
        Rosenblatt (1957) es el hito inicial, pero sus limitaciones como
        clasificador lineal estancan el área.
\end{itemize}

El redescubrimiento del algoritmo de \textbf{retropropagación del error}
en 1986 (artículo de Rumelhart, Hinton y Williams en \emph{Nature})
rehabilita los modelos conexionistas y permite el ajuste de redes
neuronales multicapa para problemas no lineales complejos, como el
reconocimiento de dígitos manuscritos. En las décadas de 1980--1990
surgen nuevos paradigmas: \textbf{sistemas fuzzy}, \textbf{máquinas de
vectores de soporte (SVM)}, \textbf{algoritmos genéticos}. Su integración
se conoció como \emph{soft computing} o \emph{inteligencia
computacional}, hoy englobando también los métodos clásicos de
estadística y \emph{machine learning}. A partir de 2012, el éxito de las
\textbf{redes neuronales profundas} (\emph{deep learning}), ejecutadas
sobre GPU, alcanza o supera el desempeño humano en tareas de visión
computacional y procesamiento de lenguaje natural.

En paralelo, la explosión de datos digitales (Internet, dispositivos
móviles desde 2007, sensores e \emph{Internet of Things}) dio origen al
término \emph{big data} (2010) y consolidó el ecosistema actual de la
ciencia de datos, dominado por lenguajes abiertos (Python, R) y
bibliotecas libres (scikit-learn, Pandas, TensorFlow).

\section{Ciencia de datos y proceso de descubrimiento de conocimiento}

En el método científico clásico, hipótesis teóricas dan origen a modelos
verificados por la observación. La ciencia de datos \emph{invierte
parcialmente} ese flujo: el objetivo es inducir modelos directamente de
los datos, en general recolectados sin planificación experimental, y por
ende sujetos a ruido, incertidumbre y sesgo muestral. Un modelo es
siempre una representación \emph{indirecta} del proceso real, válida sólo
en el dominio de la muestra que lo generó.

\begin{conceito}
\textbf{KDD — Knowledge Discovery in Databases}: proceso cíclico que va
desde la selección de datos hasta la interpretación de los modelos,
pasando por limpieza, transformación, modelado y validación. Si los
objetivos no se alcanzan, se vuelve a etapas anteriores con nuevos datos
o herramientas.
\end{conceito}

Las etapas típicas del proceso son:
\begin{enumerate}[label=\textbf{(\arabic*)}, leftmargin=1.2cm]
  \item \textbf{Selección de datos}: elección del subconjunto y de las
        variables relevantes (con participación de especialistas del
        dominio).
  \item \textbf{Pre-procesamiento}: limpieza, tratamiento de valores
        faltantes, detección de \emph{outliers}, estandarización y
        transformaciones lineales (Capítulo 2).
  \item \textbf{Modelado}: ajuste de modelos predictivos o descriptivos
        (Capítulos 3--9).
  \item \textbf{Evaluación}: estadísticas de validación, validación
        cruzada, comparación entre modelos.
  \item \textbf{Interpretación e implementación}: comunicación de los
        resultados e integración a la plataforma operacional.
\end{enumerate}

Los términos \emph{minería de datos} y \emph{ciencia de datos} se usan
con frecuencia como sinónimos. El autor sugiere la siguiente distinción:
\emph{ciencia de datos} se refiere al cuerpo de conocimiento necesario;
\emph{minería de datos} al proceso concreto de construir modelos a partir
de datos.

\section{Modelos: definiciones y taxonomía}

Un \textbf{modelo} es, en este libro, una representación matemática
simplificada de un sistema o proceso real, obtenida a partir de los datos
de una muestra. Según la naturaleza del conocimiento empleado se
distinguen:

\begin{itemize}[leftmargin=1.2cm]
  \item \textbf{Modelos de conocimiento (analíticos)}: derivados de las
        leyes fundamentales de la física o del conocimiento humano
        formalizado.
  \item \textbf{Modelos de comportamiento (entrada/salida)}: inducidos
        directamente a partir de observaciones; es el caso típico de la
        ciencia de datos.
\end{itemize}

Según el objetivo, las aplicaciones se dividen en \textbf{predictivas}
(con variables de salida a estimar) y \textbf{descriptivas} (sin
variable-objetivo, se busca una representación de los datos). La
taxonomía de algoritmos de aprendizaje es igualmente fundamental:

\begin{center}\small
\begin{tabular}{@{}l p{8.6cm}@{}}
\toprule
\textbf{Paradigma} & \textbf{Idea} \\
\midrule
Supervisado    & Pares $(\bm{x},y)$ disponibles; se ajusta
$f:\bm{x}\mapsto y$. Incluye regresión y clasificación. \\
No supervisado & Sólo $\bm{x}$; se busca estructura (grupos, reglas de
asociación, representación latente). \\
Semi-supervisado & Pequeña fracción de los datos rotulada; se usa la
parte no rotulada para mejorar el modelo. \\
Por refuerzo  & El modelo (agente) interactúa con un ambiente y se ajusta
mediante señales de recompensa. \\
\bottomrule
\end{tabular}
\end{center}

\subsection{Regresión}

Dadas variables de entrada $\bm{x}\in\R^p$ y una salida numérica
$y\in\R$, la regresión busca una función
\[
\hat{y} = f(\bm{x};\bm{\theta})
\]
cuyo \textbf{error de predicción} $e=y-\hat{y}$ se minimice en algún
sentido (típicamente cuadrático). En la regresión lineal,
$f(\bm{x})=\bm{\theta}\T\bm{x}$, eventualmente con regresores no lineales
(\emph{features}) de las entradas. La precisión del modelo se evalúa en
datos de prueba, distintos del conjunto de entrenamiento. Estos tópicos
se desarrollan en los Capítulos~3 (regresión lineal) y~4 (series
temporales y sistemas dinámicos).

\subsection{Clasificación}

En la clasificación, o \textbf{reconocimiento de patrones}, la salida es
una clase $y\in\{C_1,\dots,C_K\}$, generalmente discreta y no ordenada.
El clasificador puede verse como una búsqueda de \textbf{funciones
discriminantes}, lineales o no, que dividen el espacio de entrada en
regiones asociadas a cada clase. Se distingue entre clases
\emph{linealmente separables}, \emph{separables no linealmente} y
\emph{no separables}; en este último caso, el error mínimo (\emph{tasa
de Bayes}) es estrictamente positivo. El ejemplo histórico es el conjunto
Iris de Fisher (1936), con cuatro variables numéricas y tres clases.

\subsection{Análisis de agrupamientos (\emph{clustering})}

Sin variable de salida, se busca particionar los datos en grupos de
registros similares. El gran desafío es la \textbf{validación}, pues sin
rótulos no hay error de clasificación a minimizar; se usan índices que
reflejan propiedades geométricas (compacidad intra-cluster, separación
inter-cluster) cuya interpretación depende fuertemente del analista y
del especialista del dominio. El número de grupos suele ser parámetro de
entrada del algoritmo. Las principales abordajes (partición,
jerárquicas, espectral) se tratan en el Capítulo~6.

\subsection{Otros problemas y técnicas}

\begin{itemize}[leftmargin=1.2cm]
  \item \textbf{Reglas de asociación}: encuentran co-ocurrencias
        frecuentes entre ítems (canasta de mercado, sistemas de
        recomendación).
  \item \textbf{Aprendizaje de la representación}: nuevos métodos no
        supervisados que producen representaciones latentes útiles para
        otras tareas (Capítulo~9).
  \item \textbf{\emph{Ensembles}}: combinación de múltiples modelos para
        elevar el desempeño predictivo, técnica ganadora en competencias
        como el Premio Netflix.
\end{itemize}

\section{Recursos: bibliografía, repositorios y software}

\paragraph{Referencias teóricas.} Hastie, Tibshirani \& Friedman
(\emph{Elements of Statistical Learning}); Bishop (\emph{Pattern
Recognition and Machine Learning}); Duda, Hart \& Stork (\emph{Pattern
Classification}); Theodoridis \& Koutroumbas; Webb. Para temas
específicos: Haykin (redes neuronales), Pedrycz \& Gomide (\emph{fuzzy}),
Schölkopf \& Smola (SVM y métodos de núcleo), Vapnik (teoría del
aprendizaje estadístico), Goodfellow, Bengio \& Courville (\emph{Deep
Learning}, 2016). Para minería de datos: Han \& Kamber; Tan, Steinbach \&
Kumar; Aggarwal; Zaki \& Meira.

\paragraph{Recursos en línea.} Cursos abiertos (Coursera, Stanford, MIT;
notablemente el curso de \emph{Machine Learning} de Andrew Ng); el portal
KDnuggets (P.~Piatetsky-Shapiro); repositorios de datos como el
\emph{UCI Machine Learning Repository} y la plataforma de competencias
Kaggle. Eventos de referencia: KDD, ICDM, NeurIPS, ICML; en el área de
IC, el WCCI (FUZZ-IEEE, IJCNN, IEEE CEC).

\paragraph{Software.} El ecosistema se divide en tres capas:

\begin{itemize}[leftmargin=1.2cm]
  \item \textbf{Comerciales integrados}: SAS, IBM SPSS Modeler,
        Statistica; Oracle y Microsoft SQL Server Data Mining (minería
        dentro de la base de datos); Azure ML, Amazon ML (nube).
  \item \textbf{Libres con interfaz gráfica}: Weka (Waikato, 1997, en
        Java), RapidMiner, KNIME, Orange (Python/C++).
  \item \textbf{Lenguajes abiertos con bibliotecas}: \textbf{Python}
        (NumPy, SciPy, Pandas, scikit-learn) y \textbf{R} son, según las
        encuestas de KDnuggets, los más utilizados actualmente.
        \textbf{Matlab} sigue fuerte en ingeniería.
\end{itemize}

\section{Comentarios finales}

El capítulo introductorio fija el vocabulario y el panorama de la
disciplina. Tres ideas deben retenerse para los próximos capítulos:

\begin{enumerate}[leftmargin=1.2cm]
  \item Un modelo de ciencia de datos es \textbf{inseparable de la
        muestra} que lo generó; su validez se limita al dominio
        muestreado.
  \item El proceso KDD es \textbf{iterativo}: pre-procesamiento,
        modelado y validación se retroalimentan hasta alcanzar los
        objetivos de la aplicación.
  \item La inteligencia computacional es \textbf{plural}: combina
        estadística, aprendizaje de máquina, redes neuronales, sistemas
        \emph{fuzzy} y optimización. La elección del método debe guiarse
        por el problema, los datos disponibles y la necesidad de
        interpretabilidad.
\end{enumerate}

Los capítulos siguientes desarrollan este programa: el Capítulo~2 trata
el \textbf{pre-procesamiento}; los Capítulos~3 y~4, la \textbf{regresión}
clásica y dinámica; el Capítulo~5, la \textbf{clasificación}; el
Capítulo~6, el \textbf{análisis de agrupamientos}; y los Capítulos 7--9
los modelos no lineales de inteligencia computacional propiamente dichos.

\end{document}
